\documentclass{article}
\usepackage{amsmath}
\usepackage{amssymb}
\usepackage{amsthm}
\usepackage{MnSymbol}
\usepackage{xcolor}
\usepackage{parskip}
\usepackage{tikz}
\usepackage{bbm}
\usepackage{hyperref}
\usetikzlibrary{trees}
\usetikzlibrary{graphs}


\theoremstyle{definition}
\newtheorem{definition}{Definition}
\newtheorem{theorem}{Theorem}
\newtheorem{lemma}{Lemma}


\newcommand{\abs}[1]{\left| #1 \right|}


\title{lecture}
\date{\today}
\begin{document}
\maketitle

\section{relaxation}

distortion: if a distance is 1, then it gets converted to 17 then 17 is the distortion?

\subsection{ordinal relaxation \(\alpha\)}

cosnider \begin{equation}
    \phi:(X,\delta)\to (Y,\delta^\prime)
\end{equation}

if \(\alpha\delta(ij)<\delta(k,l)\implies\delta^\prime(ij)<\delta^\prime(kl)\)

then \(\phi\) is an oridnal embedding w relaxation \(\alpha\)

i.e. distances that are different enough should be preserved (distances that are close to being the same can be violated)


\section{tradeoffs \(\alpha\) vs d}

for every int d, n, families of metric spaces cannot be embedded unless high relaxation \(\alpha\) is allowed. i.e. \(\exists \) metric space \(T\) on n points such that the \emph{triplet relaxation } of any ordinal embedding of \(T\) into \(d\)-dimensional euclidean space will be at least \begin{equation}
    \frac{\log n}{\log d + \log\log n+8}-1
\end{equation}


\begin{definition}[girth]
    size of the smallest cycle in a agraph
\end{definition}


the idea:
create a large number of metric spaces using one graph with large girth. the number of metrics here will be greater than the number of ordinal embeddings in \(\mathbb R^a\), thus cutting just one will cause the distance between some two points to be at least the girth minus one, which means the distortion/relaxation has to be at least \(G-1\). 

key question: How dense can a graph be with a high girth?


\begin{proof}
    pick a high girth graph \(G(V,E)\) with \(\lvert V\rvert = n\).\begin{equation}
        m=\lvert E\rvert \geq \frac{1}{4}n^{1+\frac{1}{g}}
    \end{equation}
    idea, start w random graph, delete edges that create small cycles. the number of edges needed to be removed is ``not htat many''

    pick \(g=\frac{\log n}{\log d + \log\log m + 8}\)

    the number of \(\lvert E\rvert>16nd\log n\approx n log n\)


    \subsection{idea}
    now subsample the graph. large number \(N\), \(G_1,\dots G_n\) (\(G_i\) are graphs)

    for a pair \(G_i,G_j\), \(\exists v\in V\) ST in \(G_i\) \(v w\) are connected, but in \(G_j\) they are not connected. but also another pair \(v,u\) are connected in \(G_j\) but not in \(G_i\). 


    i.e. \(E_{G_j}(v)\setminus E_{G_j}(v)\neq \null\) and \(E_{G_j}(v)\setminus E_{G_i}(v)\neq \null\) where \(E(v)\) are the edges around \(v\). 

    \(v\) is called the ``witness''. in a sense, v is a vertex where \(G_{ij}\) differ.


    \subsection{continue}
    now lets do the subsampling property.

    independently sample each edge in the original graph w probability \(\frac{1}{2}\) w \(m\) edges. this generates \(N = 2^{bm}\) graphs for \(b<\frac{1}{2}\log_2 \frac{4}{3}\). 

    \begin{lemma}
        the set of subgraphs \(G_i\) satisfy the faraway property in the idea section.
    \end{lemma}

    \begin{proof}
        simplifying assumptions: 
        \begin{itemize}
            \item the original graph (the one w high girth) is \(k\) regular.
            \item independence on \(v\) being a witness for any 3 sets of \(G_ijk\)

        \end{itemize}

        the probability of \(v \) being not a witness for \(G_i,G_j\) is \(\leq 2\left( \frac{3}{4} \right)^k \)

    \end{proof}



\end{proof}








\end{document}